\section*{Chapter \ref{ch:iv:rust} --- Rust}

\begin{solution}{iq:iv:rust:1}
\texttt{best} borrows the vector immutably and is used after \texttt{push}, which needs a mutable borrow: E0502. In C++ the push
could reallocate and leave \texttt{best} dangling (\cref{iq:iv:cpp:9}). Fix by ending the borrow before mutating: copy the
\texttt{i64} out (\texttt{let best = levels[0];}), or read it after the push by index.
\omcode{code/interviews/23-rust/rust/src/lib.rs}{8}{13}{Copy the value, then mutate.}\label{lst:iv:rust:push}
\iqlookfor{the aliasing rule, the C++ bug it prevents, and a zero-cost fix.}
\end{solution}

\begin{solution}{iq:iv:rust:2}
\texttt{let archived = fills;} moves the vector; \texttt{fills} is no longer valid (E0382). Fixes: borrow instead (\texttt{let archived =
\&fills;}), which costs nothing but ties \texttt{archived}'s life to \texttt{fills}; or clone (\texttt{fills.clone()}), which copies the
heap buffer. Use \texttt{fills.len()} before the move if only the length is needed.
\iqlookfor{move semantics, and choosing between borrowing and cloning by cost.}
\end{solution}

\begin{solution}{iq:iv:rust:3}
\texttt{201} then \texttt{200}. Each \texttt{let} creates a new binding that shadows the previous one; the inner block's shadow ends with
the block. No binding is mutable: shadowing is not mutation, and the values are not overwritten.
\iqlookfor{shadowing versus mutability and scope.}
\end{solution}

\begin{solution}{iq:iv:rust:4}
The borrow checker reasons about the vector as a whole: two \texttt{\&mut} into it at once are two mutable borrows of the same
value (E0499), even for different indices. Split the borrow with \texttt{split\_at\_mut}, which returns two disjoint mutable slices
(\cref{lst:iv:rust:split}), or use indices and do the update in two statements.
\iqlookfor{why disjointness is not visible to the checker, and \texttt{split\_at\_mut}.}
\end{solution}

\begin{solution}{iq:iv:rust:5}
\texttt{fn best<'a>(a: \&'a str, b: \&'a str) -> \&'a str}. The lifetime promises that the result is valid as long as both inputs
are, so the caller may not use it after either input is gone; the compiler needs it because the result may borrow from either
(E0106).
\iqlookfor{the lifetime annotation and its meaning for callers.}
\end{solution}

\begin{solution}{iq:iv:rust:6}
\texttt{top} is a local that is dropped when the function returns, so a reference to it would dangle: E0515. Return the value,
\texttt{-> i64}, which for an integer costs nothing; for a large value return it by value (moved) or a reference into the input.
\iqlookfor{the lifetime of locals and returning by value.}
\end{solution}

\begin{solution}{iq:iv:rust:7}
\texttt{drop ignored}, \texttt{end of main}, \texttt{drop b}, \texttt{drop first}, \texttt{drop second}, \texttt{drop a}. \texttt{let \_ =}
binds nothing, so the value is dropped at once; locals are dropped in reverse declaration order (\texttt{\_b}, \texttt{\_p},
\texttt{\_a}); a struct's fields are dropped in declaration order (\texttt{first} then \texttt{second}), whereas C++ destroys members in
reverse declaration order (\cref{iq:iv:cpp:2}).
\iqlookfor{the three drop rules, including the \texttt{let \_} trap and the contrast with C++.}
\end{solution}

\begin{solution}{iq:iv:rust:8}
The iterator holds \texttt{\&'a [u8]} and a position and returns \texttt{\&'a [u8]} slices of the buffer
(\cref{lst:iv:rust:fields}): each field lives as long as the buffer, not as long as the iterator, so fields can be kept after
the iterator is gone. The crate's tests check the fields of a sample message and empty fields.
\iqlookfor{a struct with a lifetime and returned slices tied to the buffer.}
\end{solution}

\begin{solution}{iq:iv:rust:9}
\texttt{*map.entry(sym).or\_insert(0) += v;}: the entry API finds or creates the slot once and returns a mutable reference to it, where a
\texttt{get} followed by an \texttt{insert} would hash twice. The crate's \texttt{volume\_by\_symbol} does this.
\iqlookfor{the entry API and why it avoids a second lookup.}
\end{solution}

\begin{solution}{iq:iv:rust:10}
Generic: one copy per strategy type, calls resolved at compile time and inlinable, larger code, the strategy fixed at compile
time. Trait object: one copy, an indirect call through a vtable per message (a few nanoseconds and no inlining), strategies
chosen at run time from configuration. On a hot path with few strategy types, generics; for a plug-in architecture,
\texttt{dyn}. A trait is not object-safe if a method is generic, returns \texttt{Self}, or takes \texttt{self} by value without a
\texttt{Sized} bound. The crate checks that both dispatches give the same result.
\iqlookfor{monomorphisation against vtables, the trade-off, and object safety.}
\end{solution}

\begin{solution}{iq:iv:rust:11}
\texttt{Rc} counts references non-atomically, so two threads cloning or dropping it could corrupt the count: it is not
\texttt{Send}, and \texttt{thread::spawn} requires its closure to be \texttt{Send} (E0277). Use \texttt{Arc}, whose count is atomic.
\texttt{Send}: a value may be moved to another thread. \texttt{Sync}: a shared reference may be used from several threads at once (a
type is \texttt{Sync} if \texttt{\&T} is \texttt{Send}); \texttt{Cell} is \texttt{Send} but not \texttt{Sync}, \texttt{Mutex<T>} makes a
\texttt{Send} \texttt{T} usable as \texttt{Sync}.
\iqlookfor{atomic against non-atomic counting, and the precise definitions of the two traits.}
\end{solution}

\begin{solution}{iq:iv:rust:12}
The assertion that $n$ does not exceed the slice's length, checked once before the loop, guarantees that every index
$i < n$ is in bounds, so the unchecked read never goes out of bounds; the safety comment states the invariant. It would be unsound if the check were removed, weakened to a
\texttt{debug\_assert!} (absent in release builds), or placed after the loop, or if the slice could shrink during the loop (it
cannot, being borrowed immutably). Check mechanically with Miri, which interprets the code and detects out-of-bounds reads, and
with a test that the function panics on a bad \texttt{n} (the crate has one). Often the compiler removes the bounds checks of the
safe version anyway (One Quant Book~13, chapter~9).
\iqlookfor{the invariant that makes \texttt{unsafe} sound, where it would break, and Miri.}
\end{solution}

\begin{solution}{iq:iv:rust:13}
\omcode{code/interviews/23-rust/rust/src/lib.rs}{98}{115}{A shared counter with an atomic
\texttt{fetch\_add}.}\label{lst:iv:rust:count}
\texttt{fetch\_add} is a single atomic read-modify-write, so no increment is lost whatever the ordering; \texttt{Relaxed} is enough for
a pure count read after the threads are joined (joining synchronises). A stronger ordering is needed when the counter publishes
other data: a producer that writes a message and then increments a sequence number needs \texttt{Release} on the increment and
the consumer \texttt{Acquire} on the read (One Quant Book~13, chapter~11). The crate's test counts 400\,000 increments from four
threads exactly.
\iqlookfor{atomic read-modify-write, and when ordering matters beyond the counter itself.}
\end{solution}
